% !PRINTED 106
\weberpage{106}
%<<CONT
wovon jedesmal die zweite aus der ersten abgeleitet werden kann durch Vermehrung des Arguments um eine halbe Periode.

Setzt man in diesen Gleichungen $u = 0$, so folgt
\weq{5}{\begin{aligned}
A\vartheta'_{11}\left(0, \frac{\omega}{2}\right) &= \vartheta_{01}\vartheta'_{11}, & 2B\vartheta'_{11}(0, 2\omega) &= \vartheta_{10}\vartheta'_{11}, \\
A\vartheta_{10}\left(0, \frac{\omega}{2}\right) &= \vartheta_{00}\vartheta_{10}, & B\vartheta_{01}(0, 2\omega) &= \vartheta_{00}\vartheta_{01},
\end{aligned}}
woraus durch Division, mit Benutzung der Relation
\[
\vartheta'_{11} = \pi\,\vartheta_{00}\,\vartheta_{11}\,\vartheta_{01},
% [sic] Printed \vartheta_{11}; mathematically expected \vartheta_{10}.
\]
\weq{6}{\vartheta_{00}\left(0, \frac{\omega}{2}\right)\vartheta_{01}\left(0, \frac{\omega}{2}\right) = \vartheta_{01}^2,}
\weq{7}{2\vartheta_{00}(0, 2\omega)\vartheta_{10}(0, 2\omega) = \vartheta_{10}^2,}
und wenn man in (6) $\omega$ durch $2\omega$, in (7) $\omega$ durch $\omega{:}2$ ersetzt:
\weq{8}{\vartheta_{01}(0, 2\omega)^2 = \vartheta_{00}\vartheta_{01},}
\weq{9}{\vartheta_{10}\left(0, \frac{\omega}{2}\right)^2 = 2\vartheta_{00}\vartheta_{10}.}
Nach (8) und (9) ergibt sich aus den zweiten Gleichungen (5):
\[
2A = \vartheta_{10}\left(0, \frac{\omega}{2}\right), \quad B = \vartheta_{01}(0, 2\omega),
\]
und man erhält also für die \gesperrt{Gausssche Transformation}:
\weq{10}{\begin{aligned}
\vartheta_{10}\left(0, \frac{\omega}{2}\right)\vartheta_{11}\left(u, \frac{\omega}{2}\right) &= 2\vartheta_{01}(u, \omega)\vartheta_{11}(u, \omega), \\
\vartheta_{10}\left(0, \frac{\omega}{2}\right)\vartheta_{10}\left(u, \frac{\omega}{2}\right) &= 2\vartheta_{00}(u, \omega)\vartheta_{10}(u, \omega),
\end{aligned}}
und für die \gesperrt{Landensche Transformation}:
\weq{11}{\begin{aligned}
\vartheta_{01}(0, 2\omega)\vartheta_{11}(2u, 2\omega) &= \vartheta_{10}(u, \omega)\vartheta_{11}(u, \omega), \\
\vartheta_{01}(0, 2\omega)\vartheta_{01}(2u, 2\omega) &= \vartheta_{00}(u, \omega)\vartheta_{01}(u, \omega).
\end{aligned}}
Es bleiben für jede der beiden Transformationen noch zwei $\vartheta$-Funktionen auszudrücken. Man kann diese Ausdrücke aus (10), (11) herleiten nach § 21, (13), gelangt aber auch direkt dazu auf folgende Weise. Die Funktionen
\weq{12}{\vartheta_{01}\left(u, \frac{\omega}{2}\right), \quad \vartheta_{10}(2u, 2\omega)}
verschwinden für
\[
w = \frac{\omega}{4}, \quad u = \frac{1}{4}.
\]
% [sic] Printed w; the argument variable required by the formulas is u.
